\section{Module A: Single-Component Adsorption}

Module A targets single-component adsorption in a rigid framework.
The current implementation uses CO2 adsorption in MOF-5/IRMOF-1 as the first benchmark case.

\subsection{Current Scope}

Module A currently covers:

\begin{enumerate}
    \item loading a benchmark JSON input file,
    \item resolving MOF-5 to IRMOF-1,
    \item resolving the DDEC CIF file,
    \item resolving CRAFTED UFF force-field files,
    \item resolving CO2 molecule definitions,
    \item generating a LAMMPS framework data file,
    \item generating a CO2 LAMMPS molecule template,
    \item generating pair coefficients and LAMMPS force-field include files,
    \item generating one GCMC input per pressure point,
    \item executing LAMMPS pressure-point simulations,
    \item parsing LAMMPS logs,
    \item calculating absolute and excess adsorption,
    \item comparing the simulation results against reference data.
\end{enumerate}

\subsection{Current Paper-Like Setup}

\begin{longtable}{p{0.36\textwidth}p{0.48\textwidth}}
\toprule
\textbf{Setting} & \textbf{Current value} \\
\midrule
Material & MOF-5, resolved as IRMOF-1 \\
Charge scheme & DDEC \\
Adsorbate & CO2 \\
Temperature & 298.15 K \\
Pressure points & 0.1, 0.5, 1, 2, 5, 10, 20, 35 bar \\
Force field & UFF framework and CRAFTED CO2 definition \\
Framework model & Rigid \\
Unit cells & 1 x 1 x 1 \\
LAMMPS units & real \\
Pair style & \texttt{lj/cut/coul/long} \\
Cutoff & 12.8 Angstrom \\
KSpace & \texttt{pppm}, accuracy \texttt{1e-4} \\
Fugacity model & Peng-Robinson via CoolProp \\
Production cycles & 500000 \\
Initialization cycles & 500000 \\
\bottomrule
\end{longtable}

\subsection{CIF to LAMMPS Data}

The CIF converter reads the framework geometry with ASE and separately parses CIF charges from \texttt{\_atom\_site\_charge}.
This is necessary because ASE reads the geometry and cell but does not reliably preserve the charge column required for LAMMPS \texttt{atom\_style full}.

The generated framework data file contains:

\begin{itemize}
    \item atom count,
    \item atom types,
    \item restricted triclinic box parameters,
    \item masses,
    \item \texttt{Atoms \# full} section with molecule ID, atom type, charge, and Cartesian coordinates.
\end{itemize}

For IRMOF-1, the framework contains 106 atoms in the primitive simulation cell.

\subsection{Molecule Templates}

The molecule converter parses CRAFTED/RASPA \texttt{.def} files for adsorbates such as CO2 and N2.
For CO2, it creates a LAMMPS molecule template with:

\begin{itemize}
    \item three atoms,
    \item two rigid bonds,
    \item global atom type IDs,
    \item partial charges from the pseudo-atom definitions.
\end{itemize}

Because the CO2 template contains bonds, LAMMPS input generation reserves additional per-atom bond and special-neighbor slots:

\begin{lstlisting}
extra/bond/per/atom 2 extra/special/per/atom 2
\end{lstlisting}

\subsection{Force-Field Conversion}

The force-field converter parses CRAFTED UFF files and writes LAMMPS pair coefficients.
It assigns atom types for framework atoms and adsorbate atoms, then applies Lorentz-Berthelot mixing for cross interactions.

The generated force-field include file is referenced by every generated LAMMPS input script.

\subsection{LAMMPS Execution}

A short smoke test can be started with:

\begin{lstlisting}[language=bash]
conda activate MOF_sim
python benchmark.py benchmark_tao2022_mof5_co2.json --run-test
\end{lstlisting}

A full pressure-grid run can be started with:

\begin{lstlisting}[language=bash]
conda activate MOF_sim
python benchmark.py benchmark_tao2022_mof5_co2.json --run-isotherm --new-run --jobs 4
\end{lstlisting}

The \texttt{--new-run} flag creates a timestamped output directory.
The \texttt{--jobs} flag controls how many pressure points are executed in parallel.

\subsection{GCMC Interpretation}

In a GCMC simulation, CO2 molecules are inserted and deleted during the simulation.
Therefore, the number of CO2 molecules in a trajectory or dump file changes over time.
This is expected behavior and is the basis for estimating adsorption from the mean molecule count.

\subsection{Visualization}

LAMMPS trajectory dumps are written for pressure-point runs when dump output is enabled in the generated GCMC input.
The dump files can be opened with the OVITO desktop application.
Typical dump location:

\begin{lstlisting}
outputs/module_A_co2_isotherm/runs/<run_id>/dumps/gcmc_10bar.lammpstrj
\end{lstlisting}

Coloring by particle type is useful to distinguish framework atom types from inserted CO2 atom types.
